\appendix
\section{Proof Details}
\label{sec:proof-details}

\subsection{Exact ReLU Antisymmetrization}
\label{app:relu-antisymmetrization}

\begin{lemma}[Per-iterate ReLU antisymmetrization]
\label{lem:step-001-per-iterate}
For the bias-free depth-two architecture of Section~\ref{sec:setup}, every
$W\in\mathbb R^{m\times n}$, $a\in\mathbb R^m$, and
$x\in\mathcal X$ satisfy
\[
 f_{a,W}(x)-f_{a,W}(-x)=\langle W^\top a,x\rangle.
\]
No regularity, boundedness, or nonzero condition on $a$ or $W$ is required.
\end{lemma}

\begin{proof}
Let $r_j=(Wx)_j$ for $j=1,\ldots,m$.  Linearity gives
$W(-x)=-Wx$, and hence
\[
\begin{aligned}
 f_{a,W}(x)-f_{a,W}(-x)
 &=\sum_{j=1}^m a_j\bigl(\sigma(r_j)-\sigma(-r_j)\bigr).
\end{aligned}
\]
For every real $r$,
\[
 \sigma(r)-\sigma(-r)
 =\max\{0,r\}-\max\{0,-r\}=r.
\]
If $r>0$, the two ReLU values are $r$ and $0$; if $r<0$, they are $0$
and $-r$; and if $r=0$, both are zero.  Therefore
\[
\sum_{j=1}^m a_j\bigl(\sigma(r_j)-\sigma(-r_j)\bigr)
=\sum_{j=1}^m a_jr_j
=a^\top Wx
=\langle W^\top a,x\rangle.
\]
The calculation remains valid when any $r_j=0$, when $a=0$, or when
$W=0$; it uses no derivative or subgradient at a ReLU kink.
\end{proof}

\begin{proposition}[Exact latter-half aggregate linearization]
\label{prop:step-001-aggregate}
For the network and aggregate defined in Section~\ref{sec:setup}, every
realized trajectory $\omega$ and every $x\in\mathcal X$ satisfy
\[
 A_\omega(x)=G_\omega(x)-G_\omega(-x)
 =\left\langle
 \sum_{t=\lceil T/2\rceil}^{T}(W^{(t)})^\top a^{(t)},x
 \right\rangle
 =\langle v_\omega,x\rangle.
\]
\end{proposition}

\begin{proof}
Let
\[
 I_T:=\{t\in\mathbb Z:\lceil T/2\rceil\leq t\leq T\}.
\]
Because $T\geq1$, this is a finite nonempty integer interval.  Subtracting
the two aggregate sums term by term gives
\[
\begin{aligned}
 A_\omega(x)
 &=\sum_{t\in I_T}f_{a^{(t)},W^{(t)}}(x)
   -\sum_{t\in I_T}f_{a^{(t)},W^{(t)}}(-x)\\
 &=\sum_{t\in I_T}
   \bigl(f_{a^{(t)},W^{(t)}}(x)
       -f_{a^{(t)},W^{(t)}}(-x)\bigr).
\end{aligned}
\]
Applying Lemma~\ref{lem:step-001-per-iterate} separately to every listed
state and using finite linearity of the inner product yields
\[
\begin{aligned}
 A_\omega(x)
 &=\sum_{t\in I_T}\langle(W^{(t)})^\top a^{(t)},x\rangle\\
 &=\left\langle\sum_{t\in I_T}(W^{(t)})^\top a^{(t)},x\right\rangle
 =\langle v_\omega,x\rangle.
\end{aligned}
\]
Each state is merely an arbitrary finite pair $(a^{(t)},W^{(t)})$, so the
argument is valid at initialization, after any update, at every kink, and for
zero parameter states.  It invokes neither the SGD recursion nor its gradient
convention.  When $T=1$, $I_T=\{1\}$, exactly the single returned
latter-half iterate is present.  Thus the equality is pathwise and has zero
residual for every realization.
\end{proof}

Consequently, for every realized trajectory and every cube point, the
antisymmetric aggregate is the identity-coordinate score
$x\mapsto\langle v_\omega,x\rangle$.  This exact equality is the sole score
interface used in the following error-transfer argument; it remains valid in
the zero-parameter and exact/noiseless regimes.

\subsection{Tie-Aware Antipodal Error Transfer}
\label{app:error-transfer}

For a score $F:\mathcal X\to\mathbb R$, define locally
\[
 e_{F,h}(z):=\mathbf 1\!\left\{
 \operatorname{sign}_{s_0}(F(z))h(z)<0\right\}.
\]
Since both labels are in $\{-1,+1\}$, this indicator is zero exactly when
$\operatorname{sign}_{s_0}(F(z))=h(z)$.

\begin{lemma}[Tie-aware antipodal pair-error comparison]
\label{lem:step-002-pair-error}
Under Assumption~\ref{assump:antipodal-oddness}, for every
$h\in\mathcal H$, realized trajectory $\omega$, and $x\in\mathcal X$,
\[
 e_{A_\omega,h}(x)
 \leq e_{G_\omega,h}(x)+e_{G_\omega,h}(-x),
 \qquad A_\omega(x)=G_\omega(x)-G_\omega(-x).
\]
\end{lemma}

\begin{proof}
Fix the stated objects and abbreviate
\[
 y=h(x),\qquad u=G_\omega(x),\qquad q=G_\omega(-x).
\]
Assumption~\ref{assump:antipodal-oddness} gives $h(-x)=-y$.  Suppose first
that both source indicators on the right vanish.  Then $u$ is correctly
classified with label $y$, and $q$ is correctly classified with label $-y$.
The four possible pairs $(s_0,y)$ give
\[
\begin{array}{c|c|c|c|c}
s_0&y&\text{condition on }u&\text{condition on }q&
  \text{consequence for }u-q\\ \hline
+1&+1&u\geq0&q<0&u-q>0\\
+1&-1&u<0&q\geq0&u-q<0\\
-1&+1&u>0&q\leq0&u-q>0\\
-1&-1&u\leq0&q>0&u-q<0.
\end{array}
\]
Thus $y(u-q)>0$ in every row.  One correctness condition may be weak
because its target is the tie label, but the correctness condition for the
opposite label is strict.  Consequently $A_\omega(x)=u-q$ is nonzero and
has label $y$, so $e_{A_\omega,h}(x)=0$.

This table also exhausts all source-score ties.  If $u=0$ is correct, then
$y=s_0$ and correctness of $q$ for $-y$ is strict.  If $q=0$ is correct,
then $-y=s_0$ and correctness of $u$ for $y=-s_0$ is strict.  A zero at a
point labeled $-s_0$ is already a source error.  If $u=q=0$, exactly one of
the opposite labels $y$ and $-y$ equals $s_0$, so exactly one source
indicator is one.

It remains to check a tie of the antisymmetric score.  If
$A_\omega(x)=0$ and $y=s_0$, its indicator is zero.  If
$A_\omega(x)=0$ and $y=-s_0$, its indicator is one, while simultaneous
source correctness is impossible because the preceding calculation would
force $yA_\omega(x)>0$.  Hence at least one source indicator is one.  In
every remaining case where a source indicator is one, the right side is at
least one and the left side, being an indicator, is at most one.  This proves
the comparison in every zero and nonzero branch.
\end{proof}

\begin{proposition}[Exact factor-two risk transfer]
\label{prop:step-002-risk-transfer}
Under Assumption~\ref{assump:antipodal-oddness} and
Lemma~\ref{lem:step-002-pair-error}, every
$\mathcal D\in\Delta(\mathcal X)$, $h\in\mathcal H$, and realized
trajectory $\omega$ satisfy
\[
 \mathcal L_{\mathcal D,h}(A_\omega)
 \leq \mathcal L_{\mathcal D,h}(G_\omega)
      +\mathcal L_{\mathcal D^-,h}(G_\omega)
 =2\mathcal L_{\mathcal D^{\mathrm{sym}},h}(G_\omega).
\]
\end{proposition}

\begin{proof}
Multiply the pointwise inequality in
Lemma~\ref{lem:step-002-pair-error} by $\mathcal D(x)$ and sum over the
finite cube:
\[
\begin{aligned}
\mathcal L_{\mathcal D,h}(A_\omega)
&=\sum_{x\in\mathcal X}\mathcal D(x)e_{A_\omega,h}(x)\\
&\leq\sum_{x\in\mathcal X}\mathcal D(x)e_{G_\omega,h}(x)
 +\sum_{x\in\mathcal X}\mathcal D(x)e_{G_\omega,h}(-x)\\
&=\mathcal L_{\mathcal D,h}(G_\omega)
 +\sum_{x\in\mathcal X}\mathcal D(x)e_{G_\omega,h}(-x).
\end{aligned}
\]
The last indicator evaluates both the score and target at $-x$.  Since
$x\mapsto-x$ is a bijection, the substitution $z=-x$ gives
\[
\begin{aligned}
\sum_{x\in\mathcal X}\mathcal D(x)e_{G_\omega,h}(-x)
&=\sum_{z\in\mathcal X}\mathcal D(-z)e_{G_\omega,h}(z)\\
&=\sum_{z\in\mathcal X}\mathcal D^-(z)e_{G_\omega,h}(z)
 =\mathcal L_{\mathcal D^-,h}(G_\omega).
\end{aligned}
\]
Finally,
\[
\begin{aligned}
2\mathcal L_{\mathcal D^{\mathrm{sym}},h}(G_\omega)
&=2\sum_{z\in\mathcal X}
 \frac{\mathcal D(z)+\mathcal D^-(z)}2e_{G_\omega,h}(z)\\
&=\mathcal L_{\mathcal D,h}(G_\omega)
 +\mathcal L_{\mathcal D^-,h}(G_\omega).
\end{aligned}
\]
No symmetry of $\mathcal D$ is assumed.  The pointwise zero-score checks in
Lemma~\ref{lem:step-002-pair-error} remain intact under the finite sum, and
the numerical factor is exactly $2$.
\end{proof}

Together, Propositions~\ref{prop:step-001-aggregate}
and~\ref{prop:step-002-risk-transfer} attach the last risk inequality to the
exact homogeneous score $x\mapsto\langle v_\omega,x\rangle$, pathwise for
every realization and with no trajectory condition.

In particular, the pointwise comparison holds for both tie labels, all four
target/tie combinations, source-score zeros, and $A_\omega(x)=0$.  Summing
under any original distribution and changing variables by the cube involution
therefore gives exactly the factor-two symmetrized risk, with no symmetry
assumption on the original distribution.

\subsection{Distributional Extraction of an Approximate Separator}
\label{app:expected-extraction}

\begin{lemma}[Legality of antipodal symmetrization]
\label{lem:step-003-sym-law}
For every $\mathcal D\in\Delta(\mathcal X)$,
$\mathcal D^{\mathrm{sym}}=(\mathcal D+\mathcal D^-)/2$ belongs to
$\Delta(\mathcal X)$.
\end{lemma}

\begin{proof}
Every mass in $\mathcal D^{\mathrm{sym}}$ is nonnegative.  The cube
involution is a bijection, so
\[
\sum_{x\in\mathcal X}\mathcal D^-(x)
=\sum_{x\in\mathcal X}\mathcal D(-x)
=\sum_{z\in\mathcal X}\mathcal D(z)=1.
\]
Consequently,
\[
\sum_{x\in\mathcal X}\mathcal D^{\mathrm{sym}}(x)
=\frac12\left(\sum_x\mathcal D(x)+\sum_x\mathcal D^-(x)\right)=1,
\]
which is exactly the probability-distribution condition.
\end{proof}

\begin{proposition}[Expected factor-two identity-coordinate transfer]
\label{prop:step-003-expected-transfer}
Under Assumptions~\ref{assump:antipodal-oddness} and
\ref{assump:universal-sgd-success}, for every
$\mathcal D\in\Delta(\mathcal X)$ and $h\in\mathcal H$,
\[
\mathbb E_{\omega\sim\mathbb Q_{\mathcal D^{\mathrm{sym}},h}}
 \left[\mathcal L_{\mathcal D,h}
 \bigl(x\mapsto\langle v_\omega,x\rangle\bigr)\right]
\leq2\varepsilon.
\]
\end{proposition}

\begin{proof}
Fix $\mathcal D$ and $h$.  Lemma~\ref{lem:step-003-sym-law} makes
$\mathcal D^{\mathrm{sym}}$ a legal input distribution.  For every trajectory
under $\mathbb Q_{\mathcal D^{\mathrm{sym}},h}$,
Proposition~\ref{prop:step-001-aggregate} gives
\[
A_\omega(x)=\langle v_\omega,x\rangle\qquad(x\in\mathcal X),
\]
so the corresponding losses under the original evaluation law are equal.
  Applying Proposition~\ref{prop:step-002-risk-transfer} with evaluation law
  $\mathcal D$ gives the corresponding pathwise inequality.  The Gaussian
  initialization coordinates and the finitely many sampled cube points are
  measurable random variables.  Under the fixed source gradient convention,
  each one-step SGD update is a Borel function of the current coordinates and
  the current sample; induction over the finite horizon therefore makes every
  coordinate of $(a^{(t)},W^{(t)})$ measurable.  Consequently $v_\omega$ is a
  finite sum of algebraic functions of measurable coordinates.  For each
  fixed $x$, the score $\langle v_\omega,x\rangle$ is measurable, the fixed
  tie-resolved error indicator is Borel, and
  \[
  \mathcal L_{\mathcal D,h}
   \bigl(x\mapsto\langle v_\omega,x\rangle\bigr)
  =\sum_{x\in\mathcal X}\mathcal D(x)
   \mathbf 1\!\left\{
    \operatorname{sign}_{s_0}(\langle v_\omega,x\rangle)h(x)<0
   \right\}
  \]
  is a finite sum of Borel random variables.  The same finite-state argument
  makes the source loss measurable.  Both losses are bounded, so taking
  expectation under this same trajectory law gives
\[
\begin{aligned}
\mathbb E_{\mathbb Q_{\mathcal D^{\mathrm{sym}},h}}
 \mathcal L_{\mathcal D,h}\bigl(x\mapsto\langle v_\omega,x\rangle\bigr)
&\leq 2\mathbb E_{\mathbb Q_{\mathcal D^{\mathrm{sym}},h}}
 \mathcal L_{\mathcal D^{\mathrm{sym}},h}(G_\omega)\\
&\leq2\varepsilon.
\end{aligned}
\]
The second inequality is Assumption~\ref{assump:universal-sgd-success}
instantiated at $(\mathcal D^{\mathrm{sym}},h)$.  Both losses are bounded in
$[0,1]$ and the displayed pathwise inequality uses exactly the same law on
both sides; no expectation or distribution is interchanged.
\end{proof}

\begin{lemma}[Finite-domain expectation-to-existence]
\label{lem:step-003-finite-extraction}
Fix $\mathcal D\in\Delta(\mathcal X)$ and $h\in\mathcal H$.  Let
$\mathbb Q$ be a probability law for which $v_\omega$ and
$\mathcal L_{\mathcal D,h}(x\mapsto\langle v_\omega,x\rangle)$ are
measurable.  If $c\geq0$ and
\[
\mathbb E_{\omega\sim\mathbb Q}
 \left[\mathcal L_{\mathcal D,h}(x\mapsto\langle v_\omega,x\rangle)\right]
\leq c,
\]
then there is a deterministic $v\in\mathbb R^n$ with
\[
\mathcal L_{\mathcal D,h}(x\mapsto\langle v,x\rangle)\leq c.
\]
\end{lemma}

\begin{proof}
Define the finite subset-sum set
\[
 \mathscr R_{\mathcal D}:=
 \left\{\sum_{x\in\mathcal X}\mathcal D(x)b_x:
 (b_x)_{x\in\mathcal X}\in\{0,1\}^{\mathcal X}\right\}.
\]
It has at most $2^{|\mathcal X|}$ elements.  The random variable
\[
Y(\omega):=\mathcal L_{\mathcal D,h}
 \bigl(x\mapsto\langle v_\omega,x\rangle\bigr)
\]
is measurable by hypothesis and takes values in $\mathscr R_{\mathcal D}$,
because each finite-domain error indicator is either zero or one, including
at a score tie.  The events $\{Y=r\}$ partition the trajectory space, so
\[
\mathbb E_{\mathbb Q}Y
=\sum_{r\in\mathscr R_{\mathcal D}}r\,\mathbb Q[Y=r],
\qquad
\sum_{r\in\mathscr R_{\mathcal D}}\mathbb Q[Y=r]=1.
\]
If no realized trajectory had $Y(\omega)\leq c$, every value with positive
probability would be strictly larger than $c$.  The finite nonempty set of
such values would have a minimum $\gamma>c$, forcing
$\mathbb E_{\mathbb Q}Y\geq\gamma>c$, a contradiction.  Hence some value
$r\leq c$ has positive probability.  Its event is nonempty; choose a
trajectory $\omega_0$ in it and set $v=v_{\omega_0}$.

When $c=0$, nonnegativity forces a positive-probability value equal to zero,
so the selected vector has zero risk.  When $c\geq1$, every risk is already
at most $1\leq c$.  Thus no atom, margin, or limit argument is needed.
\end{proof}

\begin{proposition}[Distribution-wise approximate homogeneous separator]
\label{prop:step-003-approximate-separator}
Under Assumptions~\ref{assump:antipodal-oddness} and
\ref{assump:universal-sgd-success}, for every
$\mathcal D\in\Delta(\mathcal X)$ and $h\in\mathcal H$ there exists a
deterministic $v=v(\mathcal D,h)\in\mathbb R^n$ such that
\[
\mathcal L_{\mathcal D,h}(x\mapsto\langle v,x\rangle)\leq2\varepsilon.
\]
\end{proposition}

\begin{proof}
For the finite SGD recursion under the fixed source convention, the
measurability argument in the proof of
Proposition~\ref{prop:step-003-expected-transfer} shows that $v_\omega$ and
the displayed finite-domain tie-resolved loss are measurable under
$\mathbb Q_{\mathcal D^{\mathrm{sym}},h}$.  Apply that proposition and then
Lemma~\ref{lem:step-003-finite-extraction} with
$\mathbb Q=\mathbb Q_{\mathcal D^{\mathrm{sym}},h}$ and $c=2\varepsilon$.
The resulting vector may depend on the displayed pair $(\mathcal D,h)$, but
the identity coordinates are fixed and no common trajectory is asserted.
\end{proof}

Thus the quantifier order is
$\forall\mathcal D\,\forall h\,\exists v(\mathcal D,h)$, and the learner law
used for that pair is exactly
$\mathbb Q_{\mathcal D^{\mathrm{sym}},h}$.  The architecture, stepsize, and
finite horizon remain fixed.  When $\varepsilon=0$, the extracted vector has
zero risk; when $\mathcal H=\varnothing$, the target assertion is vacuous.

\subsection{Antipodal Representatives and Strict Feasibility}
\label{app:strict-interface}

For each fixed target $h\in\mathcal H$, define the appendix-local set
\[
 Q_h:=\{q\in\mathcal X:h(q)=-s_0\}.
\]
The next lemma proves the representative property used by the subsequent
strict-system and convex-obstruction results.

\begin{lemma}[Antipodal representative selection]
\label{lem:step-004-representatives}
Under Assumption~\ref{assump:antipodal-oddness}, for every fixed
$h\in\mathcal H$ the set $Q_h=\{q:h(q)=-s_0\}$ contains exactly one
member of each pair $\{x,-x\}$, and
\[
\mathcal X=\bigsqcup_{q\in Q_h}\{q,-q\}.
\]
\end{lemma}

\begin{proof}
For $n\geq1$, no cube point equals its antipode: $x=-x$ would force every
coordinate to be zero, contrary to $x_i\in\{-1,+1\}$.  Fix $x$ and write
$y=h(x)$.  Oddness gives $h(-x)=-y$.  Exactly one of the two opposite labels
$y$ and $-y$ equals $-s_0$, so exactly one of $x,-x$ lies in $Q_h$.  The
two-element orbits are disjoint and cover the cube, proving the partition;
for $n=1$ this is the single orbit $\{(+1),(-1)\}$.
\end{proof}

\begin{lemma}[Pairwise strict-sign criterion]
\label{lem:step-004-pairwise-strictness}
Under Assumption~\ref{assump:antipodal-oddness}, if $q\in Q_h$ and
$w\in\mathbb R^n$, then
\[
\left[
\operatorname{sign}_{s_0}(\langle w,q\rangle)=h(q)\ \text{and}\ 
\operatorname{sign}_{s_0}(\langle w,-q\rangle)=h(-q)
\right]
\Longleftrightarrow h(q)\langle w,q\rangle>0.
\]
\end{lemma}

\begin{proof}
Set $z=\langle w,q\rangle$.  Homogeneity gives
$\langle w,-q\rangle=-z$.  Since $q\in Q_h$, $h(q)=-s_0$; oddness gives
$h(-q)=s_0$.  If $s_0=+1$, then $h(q)=-1$ and
$\operatorname{sign}_{+1}(z)=h(q)$ iff $z<0$, which is iff
$h(q)z>0$.  If $s_0=-1$, then $h(q)=+1$ and
$\operatorname{sign}_{-1}(z)=h(q)$ iff $z>0$, again iff $h(q)z>0$.  Thus
\[
\operatorname{sign}_{s_0}(z)=h(q)
\Longleftrightarrow h(q)z>0.
\]
In particular $z=0$ is an error at $q$.  The forward implication follows
immediately.  Conversely, if $h(q)z>0$, then $z\neq0$ and
\[
 h(-q)\langle w,-q\rangle=(-h(q))(-z)=h(q)z>0,
\]
so both nonzero scores have their correct labels.  This proves the
equivalence, including both tie labels and the zero-score boundary.
\end{proof}

\begin{proposition}[Strict-system representation equivalence]
\label{prop:step-004-strict-equivalence}
Under Assumption~\ref{assump:antipodal-oddness}, for every fixed
$h\in\mathcal H$,
\[
\begin{aligned}
&\exists w\in\mathbb R^n\ \forall x\in\mathcal X,
\quad\operatorname{sign}_{s_0}(\langle w,x\rangle)=h(x)\\
&\qquad\Longleftrightarrow
\exists w\in\mathbb R^n\ \forall q\in Q_h,
\quad h(q)\langle w,q\rangle>0.
\end{aligned}
\]
\end{proposition}

\begin{proof}
By Lemma~\ref{lem:step-004-representatives}, the pairs indexed by $Q_h$
partition $\mathcal X$.  Exactness on all of $\mathcal X$ implies exactness
on both members of each pair, and Lemma~\ref{lem:step-004-pairwise-strictness}
then gives the strict inequality on its representative.  Conversely, the
strict inequalities give exactness on each pair by the same lemma, and the
partition extends it to every cube point.
\end{proof}

\subsection{Finite Convex Obstruction}
\label{app:convex-obstruction}

Fix $h\in\mathcal H$ and write
\[
 Z_h:=\{h(q)q:q\in Q_h\}\subseteq\mathbb R^n.
\]
By Lemma~\ref{lem:step-004-representatives}, $Q_h$ is finite and nonempty.
The strict-system equivalence in Proposition~\ref{prop:step-004-strict-equivalence}
identifies infeasibility of exact representation with the absence of a vector
having positive inner product with every member of $Z_h$.

\begin{proposition}[Finite closest-point zero certificate]
\label{prop:step-005-closest-point}
Under Assumption~\ref{assump:antipodal-oddness}, fix $h\in\mathcal H$ and set
\[
 Q_h:=\{q\in\mathcal X:h(q)=-s_0\},\qquad
 Z_h:=\{h(q)q:q\in Q_h\}.
\]
Lemma~\ref{lem:step-004-representatives} makes $Q_h$ finite and nonempty.  If
there is no $w\in\mathbb R^n$ such that
$\langle w,z\rangle>0$ for every $z\in Z_h$, then there are coefficients
$\lambda_q\geq0$ with $\sum_{q\in Q_h}\lambda_q=1$ and
\[
\sum_{q\in Q_h}\lambda_q h(q)q=0.
\]
\end{proposition}

\begin{proof}
Enumerate $Q_h=\{q_1^0,\ldots,q_r^0\}$ and put
$z_j=h(q_j^0)q_j^0$.  Every $z_j$ is nonzero and has norm $\sqrt n$.  Let
\[
 \Delta_r=\{\lambda\in\mathbb R^r:\lambda_j\geq0,
 \ \sum_{j=1}^r\lambda_j=1\},
 \qquad
 C=\left\{\sum_{j=1}^r\lambda_jz_j:\lambda\in\Delta_r\right\}.
\]
The simplex is nonempty and compact, and the continuous function
$\lambda\mapsto\|\sum_j\lambda_jz_j\|_2$ attains a minimum.  Let
$\lambda^\star$ be a minimizer and
$p=\sum_j\lambda_j^\star z_j$.

For fixed $j$ and $t\in(0,1]$, the point
$p_t=(1-t)p+t z_j$ belongs to $C$, since its coefficients are
$(1-t)\lambda_\ell^\star+t\mathbf 1\{\ell=j\}$.  Minimality of $p$ gives
\[
0\leq\|p_t\|_2^2-\|p\|_2^2
=2t\langle p,z_j-p\rangle+t^2\|z_j-p\|_2^2.
\]
After division by $t$ and passage to $t\downarrow0$,
\[
\langle p,z_j\rangle\geq\|p\|_2^2
\qquad(1\leq j\leq r).
\]
If $p\neq0$, then $w=p$ has
$\langle w,z_j\rangle\geq\|p\|_2^2>0$ for every $j$, contradicting the
local infeasibility hypothesis.  Hence $p=0$, and the minimizing coefficients
give the asserted nonnegative unit-sum representation.
\end{proof}

\begin{lemma}[Positive support pruning]
\label{lem:step-005-support-pruning}
Suppose $u_1,\ldots,u_R\in\mathbb R^n$ and $\rho_j\geq0$ satisfy
\[
\sum_{j=1}^R\rho_j=1,
\qquad
\sum_{j=1}^R\rho_j u_j=0.
\]
Then there are distinct vectors $\widetilde u_1,\ldots,\widetilde u_k$
from the original list and coefficients $\alpha_i>0$ such that
\[
1\leq k\leq n+1,
\qquad \sum_{i=1}^k\alpha_i=1,
\qquad \sum_{i=1}^k\alpha_i\widetilde u_i=0.
\]
If, for a fixed $h\in\mathcal H$ and
$Q_h:=\{q\in\mathcal X:h(q)=-s_0\}$, each retained input satisfies
$u_j=h(q_j)q_j$ for some $q_j\in Q_h$, one may retain one such $q_j$ for
each distinct $\widetilde u_i$.
\end{lemma}

\begin{proof}
Delete all indices with $\rho_j=0$; at least one coefficient remains.
If vector values are duplicated, group equal copies and replace their
coefficients by the grouped sums.  The resulting distinct list
$u_1,\ldots,u_{k_0}$ has coefficients $\beta_i>0$ with
\[
\sum_i\beta_i=1,\qquad\sum_i\beta_i u_i=0.
\]
If $k_0\leq n+1$, this is the desired representation.  Otherwise the
augmented vectors $(u_i,1)\in\mathbb R^{n+1}$ are linearly dependent, so a
nonzero $\gamma\in\mathbb R^{k_0}$ satisfies
\[
\sum_i\gamma_i u_i=0,\qquad\sum_i\gamma_i=0.
\]
Because the coordinates of $\gamma$ sum to zero, there are both positive and
negative coordinates.  Let $I_+=\{i:\gamma_i>0\}$ and define
\[
 \tau=\min_{i\in I_+}\frac{\beta_i}{\gamma_i}>0,
 \qquad
 \beta_i'=\beta_i-\tau\gamma_i.
\]
For $i\in I_+$, $\beta_i'\geq0$ and at least one is zero.  For
$\gamma_i<0$, $\beta_i'>0$, and for $\gamma_i=0$, $\beta_i'=\beta_i>0$.
Moreover,
\[
\sum_i\beta_i'=1,\qquad
\sum_i\beta_i'u_i=0.
\]
Delete the newly zero coefficients and repeat.  Each repetition strictly
reduces the finite support, so the process terminates with at most $n+1$
positive coefficients and preserves both displayed identities.  No new vector
is introduced.  A formal singleton support is allowed by the algebra; in the
present application it cannot occur because every $h(q)q$ has norm
$\sqrt n>0$, but the lower-bound argument below remains valid for $k=1$.
\end{proof}

\begin{lemma}[Signed cancellation]
\label{lem:step-005-signed-cancellation}
If $z_1,\ldots,z_k\in\mathbb R^n$ and
\[
\alpha_i>0,\qquad \sum_{i=1}^k\alpha_i=1,\qquad
\sum_{i=1}^k\alpha_i z_i=0,
\]
then for every $w\in\mathbb R^n$ there is an $i$ with
$\langle w,z_i\rangle\leq0$.
\end{lemma}

\begin{proof}
Taking the inner product with $w$ gives
\[
0=\left\langle w,\sum_{i=1}^k\alpha_i z_i\right\rangle
=\sum_{i=1}^k\alpha_i\langle w,z_i\rangle.
\]
If every inner product were strictly positive, positivity of every
$\alpha_i$ would make the finite sum strictly positive, a contradiction.
Thus one inner product is nonpositive.  This includes $w=0$ and equality
cases without replacing $\leq0$ by a strict inequality.
\end{proof}

\begin{proposition}[Uniform tie-aware obstruction]
\label{prop:step-005-uniform-obstruction}
Under Assumption~\ref{assump:antipodal-oddness}, fix $h\in\mathcal H$, set
$Q_h:=\{q\in\mathcal X:h(q)=-s_0\}$, and use the representative conclusion
of Lemma~\ref{lem:step-004-representatives}.  Suppose distinct
$q_1,\ldots,q_k\in Q_h$ and positive coefficients satisfy
\[
1\leq k\leq n+1,\qquad \sum_{i=1}^k\alpha_i=1,\qquad
\sum_{i=1}^k\alpha_i h(q_i)q_i=0.
\]
Define
\[
\mathcal D_h^\star(x)=
\begin{cases}
1/k,&x\in\{q_1,\ldots,q_k\},\\
0,&\text{otherwise}.
\end{cases}
\]
Then $\mathcal D_h^\star\in\Delta(\mathcal X)$ and, for every
$w\in\mathbb R^n$,
\[
\mathcal L_{\mathcal D_h^\star,h}(x\mapsto\langle w,x\rangle)
\geq\frac1k\geq\frac1{n+1}.
\]
\end{proposition}

\begin{proof}
Distinctness makes the displayed masses nonnegative and summing to one.  Put
$z_i=h(q_i)q_i$.  Lemma~\ref{lem:step-005-signed-cancellation} gives an
index with $\langle w,z_i\rangle\leq0$, or equivalently
\[
h(q_i)\langle w,q_i\rangle\leq0.
\]
If this product is negative, the nonzero score has prediction $-h(q_i)$ and
is an error.  If it is zero, then $q_i\in Q_h$ gives
$h(q_i)=-s_0$, while the zero score receives label $s_0=-h(q_i)$, again an
error.  This explicitly handles both choices of $s_0$.
Therefore at least one of the $k$ uniform atoms is an error, and
\[
\mathcal L_{\mathcal D_h^\star,h}(x\mapsto\langle w,x\rangle)
=\frac1k\sum_{i=1}^k
\mathbf 1\!\left\{
\operatorname{sign}_{s_0}(\langle w,q_i\rangle)h(q_i)<0
\right\}
\geq\frac1k.
\]
Since $k\leq n+1$, $1/k\geq1/(n+1)$.  For $w=0$ every atom is a tie and
the actual error is one.  The same calculation covers $k=1$, $k=n+1$, and
$n=1$.
\end{proof}

\begin{proposition}[Finite obstruction under strict-system infeasibility]
\label{prop:step-005-obstruction}
Under Assumption~\ref{assump:antipodal-oddness}, fix $h\in\mathcal H$ and set
\[
Q_h:=\{q\in\mathcal X:h(q)=-s_0\},\qquad
Z_h:=\{h(q)q:q\in Q_h\}.
\]
Use the representative conclusion of
Lemma~\ref{lem:step-004-representatives}.  If there is no
$w\in\mathbb R^n$ satisfying
\[
h(q)\langle w,q\rangle>0\qquad\text{for every }q\in Q_h
\]
(equivalently, no $w$ has positive inner product with every $z\in Z_h$),
then there are distinct representatives
$q_1,\ldots,q_k\in Q_h$, positive coefficients with $1\leq k\leq n+1$,
and a distribution $\mathcal D_h^\star$ such that
\[
\sum_{i=1}^k\alpha_i=1,\qquad
\sum_{i=1}^k\alpha_i h(q_i)q_i=0,
\qquad
\forall w\in\mathbb R^n,\quad
\mathcal L_{\mathcal D_h^\star,h}(x\mapsto\langle w,x\rangle)
\geq\frac1k\geq\frac1{n+1}.
\]
\end{proposition}

\begin{proof}
Proposition~\ref{prop:step-005-closest-point} gives a nonnegative unit-sum
zero representation of the signed vectors $h(q)q$.  Apply
Lemma~\ref{lem:step-005-support-pruning} to delete zero coefficients,
duplicates, and then affine dependencies while preserving the vector sum and
coefficient sum.  The resulting positive support has size at most $n+1$.
Proposition~\ref{prop:step-005-uniform-obstruction} applied to that support
gives the displayed uniform distribution and universal lower bound.  The
infeasibility condition is only the hypothesis of this proposition and is not
used as a primitive assumption elsewhere.
\end{proof}

\subsection{Exactification by the Strict Accuracy Gap}
\label{app:exactification}

\begin{proposition}[High-accuracy strict feasibility]
\label{prop:step-006-feasibility}
Under Assumptions~\ref{assump:antipodal-oddness},
\ref{assump:high-accuracy}, and~\ref{assump:universal-sgd-success}, for every
fixed $h\in\mathcal H$, let $Q_h:=\{q\in\mathcal X:h(q)=-s_0\}$.  Then the
strict system
\[
 h(q)\langle w_h,q\rangle>0\qquad(q\in Q_h)
\]
has a solution $w_h\in\mathbb R^n$.
\end{proposition}

\begin{proof}
Fix $h$.  Suppose for contradiction that the displayed strict system is
infeasible.  Proposition~\ref{prop:step-005-obstruction} then supplies a
distribution $\mathcal D_h^\star\in\Delta(\mathcal X)$ and an integer
$1\leq k\leq n+1$ such that, for every $w\in\mathbb R^n$,
\[
 \mathcal L_{\mathcal D_h^\star,h}(x\mapsto\langle w,x\rangle)
 \geq\frac1k\geq\frac1{n+1}.
\]
Apply Proposition~\ref{prop:step-003-approximate-separator} to this same
distribution and target.  It yields a deterministic $v\in\mathbb R^n$ with
\[
 \mathcal L_{\mathcal D_h^\star,h}(x\mapsto\langle v,x\rangle)
 \leq2\varepsilon.
\]
The lower bound is universal in $w$, so it applies to this same $v$.  The
identical target, distribution, score class, and tie-resolved loss therefore
give the strict chain
\[
 \mathcal L_{\mathcal D_h^\star,h}(\langle v,\cdot\rangle)
 \leq2\varepsilon
 <\frac1{n+1}
 \leq\frac1k
 \leq\mathcal L_{\mathcal D_h^\star,h}(\langle v,\cdot\rangle),
\]
contradicting Assumption~\ref{assump:high-accuracy}.  No expectation is
compared to a realized loss here: the deterministic vector was extracted
before the contradiction.  Thus the infeasibility hypothesis is false and a
strict separator $w_h$ exists.
\end{proof}

\begin{proposition}[Exact identity representation for a fixed target]
\label{prop:step-006-exact-representation}
Under Assumptions~\ref{assump:antipodal-oddness},
\ref{assump:high-accuracy}, and~\ref{assump:universal-sgd-success}, every
fixed $h\in\mathcal H$ has a vector $w_h\in\mathbb R^n$ satisfying
\[
\forall x\in\mathcal X,\qquad
\operatorname{sign}_{s_0}(\langle w_h,x\rangle)=h(x).
\]
\end{proposition}

\begin{proof}
Fix $h\in\mathcal H$ and set
$Q_h:=\{q\in\mathcal X:h(q)=-s_0\}$.  Proposition~\ref{prop:step-006-feasibility}
supplies a vector satisfying
$h(q)\langle w_h,q\rangle>0$ for every $q\in Q_h$.  The reverse implication in
Proposition~\ref{prop:step-004-strict-equivalence} extends its exact labels
from each representative to its antipode and hence to all of $\mathcal X$.
Strictness excludes zero scores on representatives, so the fixed tie rule is
preserved without any ordinary-sign replacement.
\end{proof}

\subsection{Common Identity Map and Dimension Bounds}
\label{app:dimension-closure}

\begin{proposition}[Common identity-map exact representation]
\label{prop:step-007-common-identity}
Under Assumptions~\ref{assump:antipodal-oddness},
\ref{assump:high-accuracy}, and~\ref{assump:universal-sgd-success}, the
single map $\varphi_{\mathrm{id}}(x)=x$ satisfies
\[
\forall h\in\mathcal H\ \exists w_h\in\mathbb R^n\ \forall x\in\mathcal X,
\qquad
\operatorname{sign}_{s_0}(\langle w_h,\varphi_{\mathrm{id}}(x)\rangle)=h(x).
\]
\end{proposition}

\begin{proof}
If $\mathcal H=\varnothing$, the target quantifier is vacuous and the map
remains fixed.  Otherwise fix an arbitrary $h\in\mathcal H$ and apply
Proposition~\ref{prop:step-006-exact-representation}.  Since
$\varphi_{\mathrm{id}}(x)=x$,
\[
\langle w_h,\varphi_{\mathrm{id}}(x)\rangle=\langle w_h,x\rangle
\]
for every $x$, and applying the same tie-resolved sign map gives the desired
identity.  The map contains no target, distribution, initialization, or
trajectory-dependent choice; only $w_h$ depends on the fixed target.
\end{proof}

\begin{proposition}[Probability-one point-mass representation]
\label{prop:step-007-point-mass}
Let $\mathcal P_{\mathrm{id}}=\delta_{\varphi_{\mathrm{id}}}$.  Under
Proposition~\ref{prop:step-007-common-identity}, for every
$\mathcal D\in\Delta(\mathcal X)$ and $h\in\mathcal H$,
\[
\Pr_{\varphi\sim\mathcal P_{\mathrm{id}}}\!\left[
\inf_{w\in\mathbb R^n}\Pr_{x\sim\mathcal D}\!\left[
\operatorname{sign}_{s_0}(\langle w,\varphi(x)\rangle)h(x)<0
\right]=0\right]=1.
\]
\end{proposition}

\begin{proof}
The point mass is fixed solely by the identity map.  For arbitrary
$(\mathcal D,h)$, Proposition~\ref{prop:step-007-common-identity} supplies
$w_h$ and hence
\[
\operatorname{sign}_{s_0}(\langle w_h,\varphi_{\mathrm{id}}(x)\rangle)h(x)
=h(x)^2=1
\qquad(x\in\mathcal X).
\]
The strict error event is therefore empty at $w=w_h$, so its probability is
zero.  Nonnegativity of every risk gives
\[
0\leq\inf_w\Pr_{x\sim\mathcal D}
 [\operatorname{sign}_{s_0}(\langle w,\varphi_{\mathrm{id}}(x)\rangle)h(x)<0]
\leq0.
\]
The infimum is exactly zero, and the sole map in the support of the Dirac law
belongs to the event with probability one.  If $\mathcal H=\varnothing$,
the universal target assertion is vacuous.
\end{proof}

\begin{lemma}[Deterministic representation implies confident representation]
\label{lem:step-007-deterministic-to-confident}
For every class $\mathcal H$, every deterministic-admissible feature
dimension is confident-admissible by taking a point mass on the same map.
Consequently,
\[
\operatorname{dc}^{1/2}(\mathcal H)\leq\operatorname{dc}(\mathcal H).
\]
\end{lemma}

\begin{proof}
Let $d$ be deterministic-admissible, with common map
$\varphi:\mathcal X\to\mathbb R^d$ and target-specific vectors $u_h$ such
that
\[
\operatorname{sign}_{s_0}(\langle u_h,\varphi(x)\rangle)=h(x)
\qquad(x\in\mathcal X).
\]
Choose $\delta_\varphi$ before $\mathcal D$ and $h$.  For every subsequent
pair, the inner risk is zero at $w=u_h$, and nonnegativity makes its infimum
zero; the point mass assigns probability one to that event.  Thus every
deterministic-admissible dimension is confident-admissible, and taking the
least dimensions gives the inequality.  The argument is vacuous but valid
when $\mathcal H=\varnothing$.
\end{proof}

\begin{proposition}[Exact dimension and parameter-count chain]
\label{prop:step-007-dimension-chain}
Under Proposition~\ref{prop:step-007-common-identity},
Lemma~\ref{lem:step-007-deterministic-to-confident}, and
$n,m,T\geq1$, $S=m(n+1)$,
\[
\operatorname{dc}^{1/2}(\mathcal H)
\leq\operatorname{dc}(\mathcal H)\leq n\leq S\leq TS.
\]
\end{proposition}

\begin{proof}
Lemma~\ref{lem:step-007-deterministic-to-confident} gives the first
inequality, and Proposition~\ref{prop:step-007-common-identity} makes the
identity map an admissible deterministic map into $\mathbb R^n$, so
$\operatorname{dc}(\mathcal H)\leq n$.  Since $m\geq1$,
\[
n\leq n+1\leq m(n+1)=S.
\]
Since $T\geq1$ and $S>0$,
\[
TS-S=(T-1)S\geq0,
\]
so $S\leq TS$.  Combining the inequalities proves the chain.  At $m=1$,
$T=1$, or $n=1$, these are respectively weak equalities or the displayed
specializations $S=n+1$, $TS=S$, and $S=2m$; no excluded boundary case is
used.  The empty-class case is covered by vacuity.
\end{proof}

\subsection{Proof of the Main Theorem}
\label{app:main-proof}

\begin{proof}
Under the three assumptions, Proposition~\ref{prop:step-006-exact-representation}
gives, for every $h\in\mathcal H$, a vector $w_h$ with exact identity-coordinate
labels.  Proposition~\ref{prop:step-007-common-identity} identifies those
coordinates with the single map $\varphi_{\mathrm{id}}$.  The point-mass
conclusion follows from Proposition~\ref{prop:step-007-point-mass}, and the
complexity and parameter inequalities follow from
Proposition~\ref{prop:step-007-dimension-chain}.  If $\mathcal H$ is empty,
all target quantifiers are vacuous and the same dimension candidate and
arithmetic remain valid.  These are exactly the assertions of
Theorem~\ref{thm:main}.
\end{proof}
